\chapter{Dónde está cada cosa}

\begin{lstlisting}
~/ros2_ws/src/
|-- mision/                    paquete principal (Python)
|   |-- mision/                nodos de ROS 2 y bibliotecas
|   |-- codigo/                programas de referencia y la red neuronal
|   |-- launch/simulacion.launch.py
|   |-- scripts/               generadores, dataset, entrenamiento, campaña, instalación
|   |-- models/                x500_busqueda, personas, árboles, suelo
|   |-- px4/                   airframe 4099_gz_x500_busqueda
|   |-- config/                GUI de Gazebo, parámetros de ArduPilot
|   |-- yolo/                  personas_yolo26s.pt (detector)
|   `-- test/                  pruebas (pytest)
|-- mision_gz_plugins/         SelectorMotores (C++), para ArduPilot con la red
|-- px4_msgs/                  mensajes de PX4 (clonado, release/1.17)
`-- ardupilot_gazebo/          plugin de ArduPilot (clonado)
\end{lstlisting}

\section{Nodos y módulos}

\begin{longtable}{>{\raggedright\arraybackslash}p{4.3cm}p{10.8cm}}
\toprule
\textbf{Archivo} & \textbf{Qué hace} \\ \midrule \endhead
\cod{detector_node.py} & YOLO sobre \cod{/camera}; publica \cod{/detecciones} y la imagen anotada. \\
\cod{localizador_node.py} & Rayo--terreno, covarianza, asociación de Mahalanobis y fusión. \\
\cod{mision_node.py} & Barrido offboard, re-observación, Neural Control, avisos GPS, evasión, métricas. \\
\cod{mision_ardupilot.py} & La misma misión con ArduCopter (GUIDED por MAVLink). \\
\cod{controlador_red.py} & La red como controlador de motores con ArduPilot. \\
\cod{geometria.py} & Proyección, localización y covarianza (compartidas con las pruebas). \\
\cod{metricas.py} & Personas localizadas, falsas alarmas, error, distancia a la ruta. \\
\cod{evasion.py}, \cod{geo.py} & Evasión con el LiDAR; conversión ENU $\rightarrow$ GPS. \\
\cod{mavproxy_trayectoria.py}, \cod{ventanas_sitl.py} & Mapa de MAVProxy con ruta y personas; acomodo de ventanas. \\
\cod{modo_neural.py}, \cod{prueba_neural.py} & Activar Neural Control a mano; escalones de prueba. \\
\bottomrule
\end{longtable}

\section{Programas de referencia (\texorpdfstring{\texttt{codigo/}}{codigo/})}

\begin{longtable}{>{\raggedright\arraybackslash}p{4.3cm}p{10.8cm}}
\toprule
\cod{sim_modelo.py} & Modelo del x500 con el control base. \\
\cod{sim_localizacion.py} & Localización, covarianza y Monte Carlo. \\
\cod{nn_imitacion.py} & Red 15-64-32-4 por clonación + DAgger (\cod{--masa}, \cod{--com}, \cod{--adelanto}). \\
\cod{ppo_refinamiento.py} & Refinamiento por PPO (\cod{--mejorado}, \cod{--forma}). \\
\cod{exportar_tflite.py} & TFLite y \cod{control_net.cpp/.hpp}. \\
\bottomrule
\end{longtable}
Sin las opciones nuevas, estos programas reproducen exactamente las cifras de referencia.

\section{Scripts}

\begin{longtable}{>{\raggedright\arraybackslash}p{4.3cm}p{10.8cm}}
\toprule
\cod{generar_mundo.py} & Bosque por oclusión, iluminación, viento y semilla, con el JSON de las personas verdaderas. \\
\cod{generar_vehiculo.py} & El x500 con cámara bajo los patines, viento y LiDAR. \\
\cod{capturar_dataset.py}, \cod{entrenar_detector.py} & Dataset con cajas visibles y entrenamiento por estrato. \\
\cod{correr_experimentos.py}, \cod{analisis_wilcoxon.py} & Campaña factorial y pruebas estadísticas. \\
\cod{instalar_entorno.sh}, \cod{instalar_px4.sh} & Instalación. \\
\bottomrule
\end{longtable}
